A star emits radiation with wavelength values in a narrow range
$\Delta\lambda \ll \lambda$, i.e. the wavelength have values between
$\lambda$ and $\lambda + \Delta\lambda$. According to Planck's relationship
(for an absolute black body), the following relation define, the energy
emitted by star in the unit of time, through the unit of area of its
surface, per length-unit of the wavelength range:
% sic: the exponent below is printed as hc/k\lambda c in the original paper
% (physically it should be hc/\lambda kT); transcribed as printed.
\[ r = \frac{2\pi h c^2}{\lambda^5\left(e^{hc/k\lambda c} - 1\right)}. \]

The spectral intensities of two radiations with wavelengths $\lambda_1$ and
respectively $\lambda_2$, both in the range $\Delta\lambda$, measured on
Earth are $I_1(\lambda_1)$ and, respectively $I_2(\lambda_2)$.
\begin{description}
\item[a.] Establish the equation which, in a general case, allows
  determining the effective temperature on the surface of the star using
  only spectral measurements.
\item[b.] Find out the approximate value of the effective temperature on the
  star surface if $hc \gg \lambda k T$.
\item[c.] Find out the relation between wavelength $\lambda_1$ and
  $\lambda_2$, if $I_1(\lambda_1) = 2 I_2(\lambda_2)$, when
  $hc \ll \lambda k T$.
\end{description}
You know: $h$ -- Planck's constant; $k$ -- Boltzmann's constant; $c$ --
speed of light in vacuum.
